% Generated by tools/build_new_dists_anthology_sections.py; do not edit by hand.

\section{Non-Haar Fourier-weighted sigma-MGFs}\label{proof:new-dist:non_haar_fourier}

\resultbox{\textbf{Canonical testing artifacts.}\par\smallskip Exact backend: \path{new_dists/non_haar_fourier_sigma_mgfs/verification.py}.\par Regression: \path{new_dists/non_haar_fourier_sigma_mgfs/test_verification.py}.\par Marimo: \path{marimo_notebooks/non_haar_fourier.py}.}

\subsection*{Scope and outcome}

For a probability law on a finite or compact matrix group, Haar averaging is
replaced by a Fourier-weighted operator rather than a Reynolds projection.
We prove the exact Schur and monomial formulas, the full matrix Fourier formula
for arbitrary random walks, its scalar character-ratio specialization for
central walks, and the Casimir-weighted heat formula.  Verification is split
by group family.  For $S_3$, every natural and deleted-representation matrix
is constructed and arbitrary convolution is compared with all irreducible
Fourier blocks.  For $U(1)$, direct heat-kernel quadrature of $2\times2$
matrices is compared with weight decomposition.  For Ewens laws on
$S_3,\ldots,S_6$, exact matrix enumeration is compared with the cycle-index
formula.  The accompanying marimo notebook exposes every comparison.

\subsection{The non-Haar master theorem}

Let $G$ be finite or compact, let $dg$ denote normalized Haar measure, and
let $\rho:G\to U(V)$ be a finite-dimensional unitary representation.  A
probability measure $\mu$ on $G$ need not be invariant.  For variables
$\mathbf t=(t_1,t_2,\ldots)$ define
\begin{equation}\label{nd:non_haar_fourier:eq:def}
 \MGF_{\mu,V}(\mathbf t)
 =\int_G\prod_{i\ge1}\det(I-t_i\rho(g))^{-1}\,d\mu(g).
\end{equation}
The statement is formal coefficientwise.  It is also analytic when only
finitely many variables are nonzero and $|t_i|<1$.

Write $p_r=\sum_i t_i^r$.  For a partition
$\tau=(\tau_1,\ldots,\tau_s)$ put
\begin{equation}\label{nd:non_haar_fourier:eq:wtau}
 W_\tau(V)=\bigotimes_{j=1}^s\Sym^{\tau_j}V.
\end{equation}

\begin{theorem}[non-Haar $\sigma$-MGF]\label{nd:non_haar_fourier:thm:master}
The following identities hold coefficientwise:
\begin{align}
 \MGF_{\mu,V}(\mathbf t)
 &=\int_G\exp\!\left(\sum_{r\ge1}\frac{p_r}{r}
     \chi_V(g^r)\right)d\mu(g),\label{nd:non_haar_fourier:eq:power}\\
 &=\sum_\lambda s_\lambda(\mathbf t)
     \int_G\chi_{S_\lambda V}(g)\,d\mu(g),\label{nd:non_haar_fourier:eq:schur}\\
 [m_\tau]\MGF_{\mu,V}
 &=\int_G\chi_{W_\tau(V)}(g)\,d\mu(g).\label{nd:non_haar_fourier:eq:monomial}
\end{align}
If
$S_\lambda V\cong\bigoplus_{\pi\in\widehat G}
m_\pi(S_\lambda V)\pi$, then
\begin{equation}\label{nd:non_haar_fourier:eq:fourier-mu}
 \boxed{\MGF_{\mu,V}
 =\sum_\lambda s_\lambda
   \sum_{\pi\in\widehat G}m_\pi(S_\lambda V)
        \Tr\widehat\mu(\pi),\qquad
 \widehat\mu(\pi)=\int_G\pi(g)\,d\mu(g).}
\end{equation}
The analogous monomial formula replaces $S_\lambda V$ by $W_\tau(V)$.
\end{theorem}

\begin{proof}
If $A$ has eigenvalues $\alpha_1,\ldots,\alpha_d$, then
\begin{equation}\label{nd:non_haar_fourier:eq:detexp}
 \det(I-tA)^{-1}
 =\prod_j(1-t\alpha_j)^{-1}
 =\exp\!\left(\sum_{r\ge1}\frac{t^r}{r}\Tr(A^r)\right).
\end{equation}
Apply this identity to $A=\rho(g)$ and multiply over $i$ to obtain
\eqref{nd:non_haar_fourier:eq:power}.  The Cauchy identity gives
\begin{equation}\label{nd:non_haar_fourier:eq:cauchy}
 \prod_i\det(I-t_i\rho(g))^{-1}
 =\sum_\lambda s_\lambda(\mathbf t)\chi_{S_\lambda V}(g),
\end{equation}
which proves \eqref{nd:non_haar_fourier:eq:schur}.  Alternatively, the coefficient of $t^q$ in
one determinant inverse is $\chi_{\Sym^qV}(g)$.  Multiplying the selected
coefficients proves \eqref{nd:non_haar_fourier:eq:monomial}.

Finally, characters add over direct sums.  On each irreducible summand $\pi$,
the integral of its character is $\Tr\int\pi(g)d\mu(g)$.  Summing with
multiplicity proves \eqref{nd:non_haar_fourier:eq:fourier-mu}.  No invariance of $\mu$ was used.
\end{proof}

\begin{remark}[what replaces Reynolds projection]
For Haar measure,
$\int_G\rho_W(g)dg$ is the idempotent projection onto $W^G$, so
$\int\chi_W=\dim W^G$.  For general $\mu$ the operator
$T_{\mu,W}=\int\rho_W(g)d\mu(g)$ is usually neither idempotent nor central.
The coefficient is its trace.  Thus a non-Haar law produces a
Fourier-weighted averaging operator, not an invariant dimension.
\end{remark}

\subsection{Random walks and the cutoff-visible spectrum}

Let $\nu$ be a step distribution and $\mu_k=\nu^{*k}$.  With the convention
in \eqref{nd:non_haar_fourier:eq:fourier-mu}, representation multiplicativity and Fubini give
\begin{equation}\label{nd:non_haar_fourier:eq:conv}
 \widehat{\nu^{*k}}(\pi)=\widehat\nu(\pi)^k.
\end{equation}

\begin{corollary}[arbitrary walk]\label{nd:non_haar_fourier:cor:walk}
For every partition $\tau$,
\begin{equation}\label{nd:non_haar_fourier:eq:walk}
 \boxed{[m_\tau]\MGF_{\nu^{*k},V}
 =\sum_{\pi\in\widehat G}m_\pi(W_\tau(V))
       \Tr\bigl(\widehat\nu(\pi)^k\bigr).}
\end{equation}
If $\nu$ is central, Schur's lemma gives
\begin{equation}\label{nd:non_haar_fourier:eq:theta}
 \widehat\nu(\pi)=\theta_\pi I_{d_\pi},\qquad
 \theta_\pi=\frac1{d_\pi}\int_G\chi_\pi(g)d\nu(g),
\end{equation}
and hence
\begin{equation}\label{nd:non_haar_fourier:eq:central}
 \boxed{[m_\tau]\MGF_{\nu^{*k},V}
 =\sum_\pi m_\pi(W_\tau(V))d_\pi\theta_\pi^k.}
\end{equation}
\end{corollary}

The trivial term in \eqref{nd:non_haar_fourier:eq:central} is the Haar coefficient.  Therefore
\begin{equation}\label{nd:non_haar_fourier:eq:deviation}
 [m_\tau](\MGF_{\nu^{*k},V}-\MGF_{\mathrm{Haar},V})
 =\sum_{\pi\ne\one}m_\pi(W_\tau(V))d_\pi\theta_\pi^k.
\end{equation}
For a fixed finite group this gives the exact leading exponential decay after
combining all terms with maximal $|\theta_\pi|$.  For a sequence
$(G_n,V_n,\nu_n)$ it also identifies the correct bounded-degree cutoff
criterion: only irreducibles occurring in the fixed $W_{n,\tau}$ are visible.
This is weaker than total-variation mixing of the full group-valued walk.

\begin{remark}[centrality is not optional]
Equation \eqref{nd:non_haar_fourier:eq:central} must not be used for a noncentral generator law.
The trace of $\widehat\nu(\pi)$ does not determine
$\Tr(\widehat\nu(\pi)^k)$.  Section~\ref{nd:non_haar_fourier:sec:s3} includes a concrete
$2\times2$ counterexample to that shortcut.
\end{remark}

\subsection{Compact heat kernels}

Let the Brownian generator be $\frac12\Delta_G$, and let $\kappa_\pi\ge0$
be the Casimir eigenvalue on $\pi$ for this normalization.  The heat
semigroup diagonalizes on matrix coefficients:
\begin{equation}\label{nd:non_haar_fourier:eq:heat-fourier}
 \widehat\mu_t(\pi)=e^{-t\kappa_\pi/2}I_{d_\pi}.
\end{equation}
Substitution in Theorem~\ref{nd:non_haar_fourier:thm:master} proves the following.

\begin{theorem}[heat formula]\label{nd:non_haar_fourier:thm:heat}
\begin{align}
 \MGF^{\mathrm{heat}}_{G,V}(t)
 &=\sum_\lambda s_\lambda
   \sum_\pi m_\pi(S_\lambda V)d_\pi e^{-t\kappa_\pi/2},
   \label{nd:non_haar_fourier:eq:heat-schur}\\
 [m_\tau]\MGF^{\mathrm{heat}}_{G,V}(t)
 &=\sum_\pi m_\pi(W_\tau(V))d_\pi e^{-t\kappa_\pi/2}.
   \label{nd:non_haar_fourier:eq:heat-monomial}
\end{align}
At $t=0$ the coefficient is $\dim W_\tau(V)$, and as $t\to\infty$ it tends
to $\dim W_\tau(V)^G$.  If $\gamma_\tau$ is the smallest positive Casimir
appearing in $W_\tau(V)$, then
\begin{equation}\label{nd:non_haar_fourier:eq:heat-asymp}
 [m_\tau]\MGF_t=\dim W_\tau(V)^G
 +e^{-t\gamma_\tau/2}
   \sum_{\kappa_\pi=\gamma_\tau}m_\pi(W_\tau)d_\pi
 +o(e^{-t\gamma_\tau/2}).
\end{equation}
\end{theorem}

This theorem applies without conceptual change to classical and exceptional
compact groups once the representation and metric normalization are fixed.
For a growing-rank statement, that normalization is essential: rescaling the
metric rescales every $\kappa_\pi$ and therefore the asserted cutoff time.

\subsection{Finite matrix group: two walks on \texorpdfstring{$S_3$}{S3}}
\label{nd:non_haar_fourier:sec:s3}

\subsubsection{Constructed representations}
For $\sigma\in S_3$, let $P_\sigma$ be its $3\times3$ permutation matrix,
$P_\sigma e_j=e_{\sigma(j)}$.  The natural module decomposes as
\begin{equation}\label{nd:non_haar_fourier:eq:natdec}
 \C^3\cong\one\oplus V_{\mathrm{std}}.
\end{equation}
The verifier also constructs the deleted matrices on
$V_{\mathrm{std}}=\{x_1+x_2+x_3=0\}$ in the basis
$f_1=e_1-e_3,f_2=e_2-e_3$.  Their entries are
\begin{equation}\label{nd:non_haar_fourier:eq:deleted}
 (D_\sigma)_{ij}=\delta_{i,\sigma(j)}-\delta_{i,\sigma(3)},
 \qquad 1\le i,j\le2,
\end{equation}
where a delta is zero if its second index is $3$.  All six matrices in each
image are generated explicitly and multiplication is exact over $\mathbb Z$.

The three irreducible character rows, ordered by identity, transposition, and
three-cycle, are
\begin{center}
\begin{tabular}{@{}lrrrr@{}}
\toprule
$\pi$ & $d_\pi$ & $\chi_\pi(1)$ & $\chi_\pi((12))$ & $\chi_\pi((123))$\\
\midrule
trivial & 1 & 1 & 1 & 1\\
sign & 1 & 1 & $-1$ & 1\\
standard & 2 & 2 & 0 & $-1$\\
\bottomrule
\end{tabular}
\end{center}

\subsubsection{Central lazy-transposition law}
Set
\begin{equation}\label{nd:non_haar_fourier:eq:centralstep}
 \nu_c=\frac14\delta_e+\frac14\sum_{r\text{ a transposition}}\delta_r.
\end{equation}
It is central, and \eqref{nd:non_haar_fourier:eq:theta} gives
\begin{equation}\label{nd:non_haar_fourier:eq:s3theta}
 \theta_{\mathrm{triv}}=1,\qquad
 \theta_{\mathrm{sign}}=-\frac12,\qquad
 \theta_{\mathrm{std}}=\frac14.
\end{equation}
Thus every tested coefficient has the completely explicit form
\begin{equation}\label{nd:non_haar_fourier:eq:s3centralformula}
 m_{\mathrm{triv}}(W_\tau)
 +m_{\mathrm{sign}}(W_\tau)(-1/2)^k
 +2m_{\mathrm{std}}(W_\tau)(1/4)^k.
\end{equation}

\subsubsection{Noncentral two-generator law}
Set
\begin{equation}\label{nd:non_haar_fourier:eq:noncentralstep}
 \nu_{nc}=\frac12\delta_{(23)}+\frac12\delta_{(123)}.
\end{equation}
Using \eqref{nd:non_haar_fourier:eq:deleted}, its standard Fourier block is
\begin{equation}\label{nd:non_haar_fourier:eq:noncentralmatrix}
 \widehat\nu_{nc}(\mathrm{std})
 =\frac12D_{(23)}+\frac12D_{(123)}
 =\begin{pmatrix}0&-1/2\\0&-1/2\end{pmatrix}.
\end{equation}
This is not scalar.  The correct term in \eqref{nd:non_haar_fourier:eq:walk} is
$m_{\mathrm{std}}(W_\tau)\Tr(\widehat\nu_{nc}(\mathrm{std})^k)$;
a character-ratio substitution would discard matrix information.

\subsubsection{Exact verification strategy}
For each representation, partition $\tau$ through total degree four, and
$k\in\{0,1,2,3,5\}$, the program compares:
\begin{enumerate}
\item \textbf{Direct convolution and matrices.}  It forms $\nu^{*k}$ on all
six elements.  For each matrix $A$, it computes
\begin{equation}\label{nd:non_haar_fourier:eq:newton}
 h_0(A)=1,\qquad h_d(A)=\frac1d\sum_{r=1}^d
 \Tr(A^r)h_{d-r}(A),
\end{equation}
then averages $\prod_jh_{\tau_j}(A)$ literally.
\item \textbf{Full Fourier route.}  Multiplicities are computed independently
from
\begin{equation}\label{nd:non_haar_fourier:eq:multiplicity}
 m_\pi(W_\tau)=\frac1{6}\sum_{g\in S_3}
 \chi_{W_\tau}(g)\chi_\pi(g^{-1}),
\end{equation}
and inserted into \eqref{nd:non_haar_fourier:eq:walk} using exact rational matrix powers.
\item \textbf{Central scalar route.}  For $\nu_c$ only, the code separately
checks that every Fourier block is scalar and evaluates
\eqref{nd:non_haar_fourier:eq:s3centralformula}.
\end{enumerate}

\begin{center}
\begin{tabular}{@{}lllrrc@{}}
\toprule
walk & representation & $(k,\tau)$ & direct & Fourier & result\\
\midrule
central & natural & $(3,(1))$ & $33/32$ & $33/32$ & \textcolor{teal}{\textbf{PASS}}\\
central & natural & $(3,(2,2))$ & $63/8$ & $63/8$ & \textcolor{teal}{\textbf{PASS}}\\
central & deleted & $(5,(1))$ & $1/512$ & $1/512$ & \textcolor{teal}{\textbf{PASS}}\\
noncentral & natural & $(3,(1))$ & $7/8$ & $7/8$ & \textcolor{teal}{\textbf{PASS}}\\
noncentral & deleted & $(3,(2,1))$ & $3/4$ & $3/4$ & \textcolor{teal}{\textbf{PASS}}\\
noncentral & deleted & $(5,(2,2))$ & $61/32$ & $61/32$ & \textcolor{teal}{\textbf{PASS}}\\
\bottomrule
\end{tabular}
\end{center}
The displayed noncentral values are generated by the same suite as the full
table.  All 220 exact comparisons pass.

\subsection{Compact matrix group: \texorpdfstring{$U(1)$}{U(1)} heat flow}

Take
\begin{equation}\label{nd:non_haar_fourier:eq:u1rep}
 \rho(e^{i\phi})=\begin{pmatrix}e^{i\phi}&0\\0&e^{-i\phi}\end{pmatrix}
\end{equation}
on a two-dimensional representation with weights $+1,-1$.  For the standard
circle Laplacian, the weight-$q$ character has heat multiplier
$e^{-tq^2/2}$.  Moreover
\begin{equation}\label{nd:non_haar_fourier:eq:symweights}
 \chi_{\Sym^dV}(e^{i\phi})
 =\sum_{a=0}^d e^{i(d-2a)\phi}.
\end{equation}
If $c_\tau(q)$ is the coefficient of $z^q$ in
\begin{equation}\label{nd:non_haar_fourier:eq:weightpoly}
 \prod_{j=1}^{\ell(\tau)}
   (z^{\tau_j}+z^{\tau_j-2}+\cdots+z^{-\tau_j}),
\end{equation}
then Theorem~\ref{nd:non_haar_fourier:thm:heat} becomes the finite exact sum
\begin{equation}\label{nd:non_haar_fourier:eq:u1formula}
 \boxed{[m_\tau]\MGF_t
 =\sum_{q\in\mathbb Z}c_\tau(q)e^{-tq^2/2}.}
\end{equation}
At $t=0$ this is $\prod_j(\tau_j+1)$; at Haar time it tends to
$c_\tau(0)$.  The first nonzero $|q|$ in \eqref{nd:non_haar_fourier:eq:weightpoly} gives the
leading decay rate, an explicit instance of \eqref{nd:non_haar_fourier:eq:heat-asymp}.

\subsubsection{Independent numerical route}
The verifier does not use \eqref{nd:non_haar_fourier:eq:weightpoly} in its direct route.  On a
256-node periodic grid it constructs \eqref{nd:non_haar_fourier:eq:u1rep}, recovers every
$h_d(\rho(e^{i\phi}))$ from matrix traces using \eqref{nd:non_haar_fourier:eq:newton}, and
multiplies by
\begin{equation}\label{nd:non_haar_fourier:eq:heatkernel}
 K_t(\phi)=1+2\sum_{q=1}^{48}e^{-tq^2/2}\cos(q\phi).
\end{equation}
The trapezoidal average is compared with \eqref{nd:non_haar_fourier:eq:u1formula}.  The grid
resolves all retained Fourier modes; the omitted tail is far below the
reported tolerance at the smallest tested time $t=0.2$.

\begin{center}
\begin{tabular}{@{}crrrc@{}}
\toprule
$t$ & $\tau$ & direct quadrature & weight formula & error\\
\midrule
$0.2$ & $(1)$ & $1.809674836071920$ & $1.809674836071919$ & $6.7\times10^{-16}$\\
$0.2$ & $(2,2)$ & $6.085073220131871$ & $6.085073220131868$ & $3.6\times10^{-15}$\\
$0.7$ & $(1,1)$ & $2.493193927883212$ & $2.493193927883213$ & $1.3\times10^{-15}$\\
$2.0$ & $(5)$ & $0.736005701978834$ & $0.736005701978834$ & $4.4\times10^{-16}$\\
\bottomrule
\end{tabular}
\end{center}
All 54 comparisons through degree five pass.  The largest absolute error over
the complete suite is $1.5988\times10^{-14}$.  For all 18 partitions, the
$t=0$ dimension and $t\to\infty$ Haar-weight identities also pass.

\subsection{Ewens laws on symmetric matrix groups}

For $\sigma\in S_n$, let $K(\sigma)$ be its number of cycles and take
\begin{equation}\label{nd:non_haar_fourier:eq:ewenslaw}
 \mathbb P_{n,\theta}(\sigma)
 =\frac{\theta^{K(\sigma)}}{\theta^{(n)}},\qquad
 \theta^{(n)}=\theta(\theta+1)\cdots(\theta+n-1).
\end{equation}
The normalization follows from
$\sum_{\sigma\in S_n}\theta^{K(\sigma)}=\theta^{(n)}$.

If $m_j(\lambda)$ is the number of parts of size $j$ in $\lambda\vdash n$
and $z_\lambda=\prod_jj^{m_j(\lambda)}m_j(\lambda)!$, a permutation of cycle
shape $\lambda$ satisfies
\begin{equation}\label{nd:non_haar_fourier:eq:cycledet}
 \det(I-tP_\sigma)^{-1}
 =\prod_{j\ge1}(1-t^j)^{-m_j(\lambda)}.
\end{equation}

\begin{proposition}[exact Ewens formula]\label{nd:non_haar_fourier:prop:ewens}
For the natural permutation representation,
\begin{equation}\label{nd:non_haar_fourier:eq:ewens}
 \boxed{\MGF_{n,\theta}(\mathbf t)
 =\frac{n!}{\theta^{(n)}}
  \sum_{\lambda\vdash n}\frac{\theta^{\ell(\lambda)}}{z_\lambda}
  \prod_{i,j}(1-t_i^j)^{-m_j(\lambda)}.}
\end{equation}
\end{proposition}

\begin{proof}
There are $n!/z_\lambda$ permutations of shape $\lambda$.  They all have
$\ell(\lambda)$ cycles, hence the same Ewens weight, and
\eqref{nd:non_haar_fourier:eq:cycledet} gives their common matrix integrand.  Grouping the direct
element sum by shape proves \eqref{nd:non_haar_fourier:eq:ewens}.
\end{proof}

For fixed $j$, the Ewens cycle counts converge jointly to independent
Poisson variables of means $\theta/j$.  Consequently, coefficientwise in the
positive-degree variables,
\begin{equation}\label{nd:non_haar_fourier:eq:ewenslimit}
 \MGF_{\theta,\infty}
 =\exp\!\left[\sum_{j\ge1}\frac{\theta}{j}
   \left(\prod_i(1-t_i^j)^{-1}-1\right)\right].
\end{equation}
For a coefficient of bounded degree, only finitely many $j$ contribute, so
the formal interpretation is unambiguous.

\subsubsection{Exact verification strategy and results}
The direct route enumerates all $n!$ permutation matrices, computes their
symmetric characters from matrix powers, and applies \eqref{nd:non_haar_fourier:eq:ewenslaw}.
The independent route stores only cycle lengths, expands
\eqref{nd:non_haar_fourier:eq:cycledet}, and aggregates class sizes.  Both use rational
arithmetic for
\begin{equation*}
 n=3,4,5,6,\qquad \theta\in\{1/2,1,2\},
 \qquad |\tau|\le4.
\end{equation*}
This gives 132 coefficient comparisons.  At
$(t_1,t_2)=(1/10,-1/20)$, 12 further comparisons evaluate the full rational
determinant average without truncating in degree.  Every comparison passes.

\begin{center}
\begin{tabular}{@{}cccrrc@{}}
\toprule
$n$ & $\theta$ & $\tau$ & matrix enumeration & cycle formula & result\\
\midrule
3 & $1/2$ & $(1)$ & $3/5$ & $3/5$ & \textcolor{teal}{\textbf{PASS}}\\
3 & $1/2$ & $(2,2)$ & $4$ & $4$ & \textcolor{teal}{\textbf{PASS}}\\
3 & $2$ & $(1)$ & $3/2$ & $3/2$ & \textcolor{teal}{\textbf{PASS}}\\
6 & $1/2$ & $(2,1)$ & $326/231$ & $326/231$ & \textcolor{teal}{\textbf{PASS}}\\
6 & $1$ & $(2,2)$ & $9$ & $9$ & \textcolor{teal}{\textbf{PASS}}\\
6 & $2$ & $(1,1)$ & $32/7$ & $32/7$ & \textcolor{teal}{\textbf{PASS}}\\
\bottomrule
\end{tabular}
\end{center}

\subsection{Related measures and representation choices}

\subsubsection{Conjugacy-class walks}
If $\nu_K$ is uniform on a conjugacy class $K$, then
$\theta_\pi=\chi_\pi(K)/d_\pi$.  Therefore
\begin{equation}\label{nd:non_haar_fourier:eq:classwalk}
 [m_\tau]\MGF_{\nu_K^{*k},V}
 =\sum_\pi m_\pi(W_\tau)d_\pi
   \left(\frac{\chi_\pi(K)}{d_\pi}\right)^k.
\end{equation}
This includes random transpositions and central reflection walks.  If a
generator set is a union of classes, its weighted average character ratio is
used.  Noncentral elementary-generator walks require \eqref{nd:non_haar_fourier:eq:walk} instead.

\subsubsection{Character-weighted laws}
For an irreducible $\eta$, character orthogonality makes
$|\chi_\eta(g)|^2dg$ a probability measure.  Formula \eqref{nd:non_haar_fourier:eq:monomial}
becomes
\begin{equation}\label{nd:non_haar_fourier:eq:charweighted}
 [m_\tau]\MGF_{|\chi_\eta|^2,V}
 =\langle\chi_{W_\tau},\chi_\eta\chi_{\eta^*}\rangle_G
 =\dim\Hom_G(W_\tau,\eta\otimes\eta^*).
\end{equation}
This is a tensor-product multiplicity, not generally an invariant dimension.
Plancherel measure itself lives on $\widehat G$, so an additional construction
is needed before it defines a law on group elements.

\subsubsection{Finite-field and Clifford cautions}
A finite-field matrix action is not automatically a complex representation.
One must specify a complex permutation, Weil, or other cross-characteristic
representation before applying the theorem.  For scalar-rich unitary groups,
one-sided moments may vanish in many degrees; a balanced series involving
both $g$ and $g^*$ can be more informative.  These are changes of
representation or observable, not exceptions to Theorem~\ref{nd:non_haar_fourier:thm:master}.

\begingroup
\small
\setlength{\parskip}{3pt}
\subsection{Reproducibility and audit summary}

The package contains one pure-Python verifier and one marimo interface.  No
external character table, symbolic algebra service, or random seed is used.
Finite computations use \texttt{fractions.Fraction}; only the compact
quadrature uses floating point.

\begin{center}
\begin{tabularx}{\linewidth}{@{}lXrrc@{}}
\toprule
family & independent routes & cases & comparisons & result\\
\midrule
$S_3$ walks & convolution/matrix traces; full Fourier blocks; central scalars
& $2$ laws, $2$ reps & 220 & \textcolor{teal}{\textbf{PASS}}\\
$U(1)$ heat & matrix-integrand quadrature; weight/Casimir sum
& $3$ times, $18$ partitions & 54 & \textcolor{teal}{\textbf{PASS}}\\
Ewens $S_n$ & element matrices; cycle-shape aggregation
& $4$ ranks, $3$ parameters & 132 & \textcolor{teal}{\textbf{PASS}}\\
Ewens full series & element determinants; class determinants
& $12$ laws & 12 & \textcolor{teal}{\textbf{PASS}}\\
\bottomrule
\end{tabularx}
\end{center}

Run the tests from the repository root:
\begin{verbatim}
.venv/bin/python -m pytest -q \
  new_dists/non_haar_fourier_sigma_mgfs/test_verification.py
\end{verbatim}
Open the grouped interactive notebook with:
\begin{verbatim}
.venv/bin/marimo edit \
  marimo_notebooks/non_haar_fourier.py
\end{verbatim}
\resultbox{\textbf{Conclusion.}  The tested cases verify the intended
replacement precisely: Haar measure extracts the trivial representation,
whereas a non-Haar law weights every constituent by a Fourier trace.  Central
walks reduce to character ratios, noncentral walks retain full Fourier
matrices, heat flow supplies Casimir exponentials, and Ewens weighting is
captured exactly by cycle-shape aggregation.}
\endgroup
